\documentclass[10pt,a4paper]{amsart}
\usepackage{amsmath,amssymb,amsthm,graphicx}
\usepackage{verbatim}
\newenvironment{reference}{\comment}{\endcomment}
%\newenvironment{reference}{}{}
\newenvironment{slogan}{\comment}{\endcomment}
\newenvironment{history}{\comment}{\endcomment}
% For commutative diagrams we use Xy-pic
\usepackage[all]{xy}

% We use 2cell for 2-commutative diagrams.
\xyoption{2cell}
% Theorem environments.
%
\theoremstyle{plain}
\newtheorem{theorem}[subsection]{Theorem}
\newtheorem{proposition}[subsection]{Proposition}
\newtheorem{lemma}[subsection]{Lemma}

\theoremstyle{definition}
\newtheorem{definition}[subsection]{Definition}
\newtheorem{example}[subsection]{Example}
\newtheorem{exercise}[subsection]{Exercise}
\newtheorem{situation}[subsection]{Situation}

\theoremstyle{remark}
\newtheorem{remark}[subsection]{Remark}
\newtheorem{remarks}[subsection]{Remarks}

\numberwithin{equation}{subsection}

% Macros
%
\def\lim{\mathop{\mathrm{lim}}\nolimits}
\def\colim{\mathop{\mathrm{colim}}\nolimits}
\def\Spec{\mathop{\mathrm{Spec}}}
\def\Hom{\mathop{\mathrm{Hom}}\nolimits}
\def\Ext{\mathop{\mathrm{Ext}}\nolimits}
\def\SheafHom{\mathop{\mathcal{H}\!\mathit{om}}\nolimits}
\def\SheafExt{\mathop{\mathcal{E}\!\mathit{xt}}\nolimits}
\def\Sch{\mathit{Sch}}
\def\Mor{\mathop{\mathrm{Mor}}\nolimits}
\def\Ob{\mathop{\mathrm{Ob}}\nolimits}
\def\Sh{\mathop{\mathit{Sh}}\nolimits}
\def\NL{\mathop{N\!L}\nolimits}
\def\CH{\mathop{\mathrm{CH}}\nolimits}
\def\proetale{{pro\text{-}\acute{e}tale}}
\def\etale{{\acute{e}tale}}
\def\QCoh{\mathit{QCoh}}
\def\Ker{\mathop{\mathrm{Ker}}}
\def\Im{\mathop{\mathrm{Im}}}
\def\Coker{\mathop{\mathrm{Coker}}}
\def\Coim{\mathop{\mathrm{Coim}}}

% Boxtimes
%
\DeclareMathSymbol{\boxtimes}{\mathbin}{AMSa}{"02}

%
% Macros for moduli stacks/spaces
%
\def\QCohstack{\mathcal{QC}\!\mathit{oh}}
\def\Cohstack{\mathcal{C}\!\mathit{oh}}
\def\Spacesstack{\mathcal{S}\!\mathit{paces}}
\def\Quotfunctor{\mathrm{Quot}}
\def\Hilbfunctor{\mathrm{Hilb}}
\def\Curvesstack{\mathcal{C}\!\mathit{urves}}
\def\Polarizedstack{\mathcal{P}\!\mathit{olarized}}
\def\Complexesstack{\mathcal{C}\!\mathit{omplexes}}
% \Pic is the operator that assigns to X its picard group, usage \Pic(X)
% \Picardstack_{X/B} denotes the Picard stack of X over B
% \Picardfunctor_{X/B} denotes the Picard functor of X over B
\def\Pic{\mathop{\mathrm{Pic}}\nolimits}
\def\Picardstack{\mathcal{P}\!\mathit{ic}}
\def\Picardfunctor{\mathrm{Pic}}
\def\Deformationcategory{\mathcal{D}\!\mathit{ef}}
\usepackage{cleveref}
\crefname{originaldefinition}{original definition}{original definitions}
\Crefname{originaldefinition}{Original definition}{Original definitions}
\makeatletter
\newlabel{definition-separated}{{definition-separated}{source}{}{sourcepointer-1}{}}
\newlabel{definition-separated@cref}{{[originaldefinition][0][]definition-separated}{[1][0][]source}{}{}{}}
\makeatother
\crefname{originallemma}{original lemma}{original lemmas}
\Crefname{originallemma}{Original lemma}{Original lemmas}
\makeatletter
\newlabel{lemma-fibre-product-after-map}{{lemma-fibre-product-after-map}{source}{}{sourcepointer-2}{}}
\newlabel{lemma-fibre-product-after-map@cref}{{[originallemma][0][]lemma-fibre-product-after-map}{[1][0][]source}{}{}{}}
\makeatother
\crefname{originalreference}{original reference}{original references}
\Crefname{originalreference}{Original reference}{Original references}
\makeatletter
\newlabel{spaces-morphisms-lemma-composition-proper}{{spaces-morphisms-lemma-composition-proper}{source}{}{sourcepointer-3}{}}
\newlabel{spaces-morphisms-lemma-composition-proper@cref}{{[originalreference][0][]spaces-morphisms-lemma-composition-proper}{[1][0][]source}{}{}{}}
\makeatother
\crefname{originalreference}{original reference}{original references}
\Crefname{originalreference}{Original reference}{Original references}
\makeatletter
\newlabel{spaces-morphisms-lemma-composition-quasi-compact}{{spaces-morphisms-lemma-composition-quasi-compact}{source}{}{sourcepointer-4}{}}
\newlabel{spaces-morphisms-lemma-composition-quasi-compact@cref}{{[originalreference][0][]spaces-morphisms-lemma-composition-quasi-compact}{[1][0][]source}{}{}{}}
\makeatother
\crefname{originalreference}{original reference}{original references}
\Crefname{originalreference}{Original reference}{Original references}
\makeatletter
\newlabel{spaces-morphisms-lemma-composition-quasi-finite}{{spaces-morphisms-lemma-composition-quasi-finite}{source}{}{sourcepointer-5}{}}
\newlabel{spaces-morphisms-lemma-composition-quasi-finite@cref}{{[originalreference][0][]spaces-morphisms-lemma-composition-quasi-finite}{[1][0][]source}{}{}{}}
\makeatother
\crefname{originalreference}{original reference}{original references}
\Crefname{originalreference}{Original reference}{Original references}
\makeatletter
\newlabel{spaces-morphisms-lemma-composition-separated}{{spaces-morphisms-lemma-composition-separated}{source}{}{sourcepointer-6}{}}
\newlabel{spaces-morphisms-lemma-composition-separated@cref}{{[originalreference][0][]spaces-morphisms-lemma-composition-separated}{[1][0][]source}{}{}{}}
\makeatother
\crefname{originalreference}{original reference}{original references}
\Crefname{originalreference}{Original reference}{Original references}
\makeatletter
\newlabel{spaces-morphisms-lemma-composition-unramified}{{spaces-morphisms-lemma-composition-unramified}{source}{}{sourcepointer-7}{}}
\newlabel{spaces-morphisms-lemma-composition-unramified@cref}{{[originalreference][0][]spaces-morphisms-lemma-composition-unramified}{[1][0][]source}{}{}{}}
\makeatother
\crefname{originalreference}{original reference}{original references}
\Crefname{originalreference}{Original reference}{Original references}
\makeatletter
\newlabel{stacks-properties-section-properties-morphisms}{{stacks-properties-section-properties-morphisms}{source}{}{sourcepointer-8}{}}
\newlabel{stacks-properties-section-properties-morphisms@cref}{{[originalreference][0][]stacks-properties-section-properties-morphisms}{[1][0][]source}{}{}{}}
\makeatother
\title{The separation axioms are stable under composition of morphisms of algebraic stacks}
\author{The Stacks Project Authors}
\date{stacks-project:stacks-morphisms.tex:local-snapshot-20261012; GNU FDL 1.2 or later}
\begin{document}
\maketitle
\noindent\textit{Editorial note.} All delivered mathematical text is the author's own, copied from the pinned Stacks Project chapter source sources/stacks-algebraic-stacks-morphisms/pinned-original-source.tex; no statement or proof step was rewritten, reordered or removed.\par\medskip
\noindent\textit{Editorial note.} Preamble provenance: the Stacks Project keeps one single shared preamble for every chapter; this card takes it directly from the staged sources/stacks-algebraic-stacks-morphisms/preamble.tex as three verbatim spans (lines 8-12, 14-18 and 154-217), so no derived file was written and nothing was added or redefined. The wrapper class and the base packages are supplied by the assembler.\par\medskip
\noindent\textit{Editorial note.} Bibliography provenance: neither the counted statement nor its proof contains a citation (this lane rescans the whole span before writing the card), so the card carries no bibliography and no entry was reconstructed.\par\medskip
\noindent\textit{Editorial note.} Selection provenance: this card was opened under the file-based control contract (CONTROL.md), with the mechanical pre-card gate run on its statement and proof interval before the card was written.\par\medskip
\noindent\textit{Allowed standard prerequisites (editorial).}\begin{itemize}
\item The factorization of the diagonal of a composite through the fibre product over the intermediate stack. Used at the start of the counted proof (source lines 859-866). Splits the diagonal into the graph of the first morphism and the base change of the second.
\item The statement that the comparison morphism between the two fibre products has the separation property. Invoked inside the counted proof (source lines 867-870). Gives the property for the second factor of the factorization.
\item The composition permanence of the corresponding properties of morphisms of algebraic spaces. Invoked at the end of the counted proof (source lines 871-884) as cross-chapter pointers. Combines the two factors into the asserted property of the composite diagonal.
\end{itemize}\medskip
\ifdefined\hypertarget\hypertarget{sourcepointer-1}{}\fi
\noindent\textit{Source reference.} Original definition definition-separated: sources/stacks-algebraic-stacks-morphisms/pinned-original-source.tex, source lines 601--601. Navigation pointer only: the counted text refers to the Stacks item definition-separated, which is an existing labeled result of this same chapter and is not part of this card. The pointer carries the target's own label as its number and its definition line as its span; its content is not reproduced here.\par\medskip
\ifdefined\hypertarget\hypertarget{sourcepointer-2}{}\fi
\noindent\textit{Source reference.} Original lemma lemma-fibre-product-after-map: sources/stacks-algebraic-stacks-morphisms/pinned-original-source.tex, source lines 763--763. Navigation pointer only: the counted text refers to the Stacks item lemma-fibre-product-after-map, which is an existing labeled result of this same chapter and is not part of this card. The pointer carries the target's own label as its number and its definition line as its span; its content is not reproduced here.\par\medskip
\ifdefined\hypertarget\hypertarget{sourcepointer-3}{}\fi
\noindent\textit{Source reference.} Original reference spaces-morphisms-lemma-composition-proper: sources/stacks-algebraic-stacks-morphisms/pinned-original-source.tex, source lines 882--882. Navigation pointer only: the counted text refers to the Stacks item spaces-morphisms-lemma-composition-proper, which belongs to another chapter of the project and is not part of this card. The pointer carries the target's own label as its number because this chapter's staged snapshot does not contain the other chapter's tag table; its content is not reproduced here.\par\medskip
\ifdefined\hypertarget\hypertarget{sourcepointer-4}{}\fi
\noindent\textit{Source reference.} Original reference spaces-morphisms-lemma-composition-quasi-compact: sources/stacks-algebraic-stacks-morphisms/pinned-original-source.tex, source lines 883--883. Navigation pointer only: the counted text refers to the Stacks item spaces-morphisms-lemma-composition-quasi-compact, which belongs to another chapter of the project and is not part of this card. The pointer carries the target's own label as its number because this chapter's staged snapshot does not contain the other chapter's tag table; its content is not reproduced here.\par\medskip
\ifdefined\hypertarget\hypertarget{sourcepointer-5}{}\fi
\noindent\textit{Source reference.} Original reference spaces-morphisms-lemma-composition-quasi-finite: sources/stacks-algebraic-stacks-morphisms/pinned-original-source.tex, source lines 881--881. Navigation pointer only: the counted text refers to the Stacks item spaces-morphisms-lemma-composition-quasi-finite, which belongs to another chapter of the project and is not part of this card. The pointer carries the target's own label as its number because this chapter's staged snapshot does not contain the other chapter's tag table; its content is not reproduced here.\par\medskip
\ifdefined\hypertarget\hypertarget{sourcepointer-6}{}\fi
\noindent\textit{Source reference.} Original reference spaces-morphisms-lemma-composition-separated: sources/stacks-algebraic-stacks-morphisms/pinned-original-source.tex, source lines 884--884. Navigation pointer only: the counted text refers to the Stacks item spaces-morphisms-lemma-composition-separated, which belongs to another chapter of the project and is not part of this card. The pointer carries the target's own label as its number because this chapter's staged snapshot does not contain the other chapter's tag table; its content is not reproduced here.\par\medskip
\ifdefined\hypertarget\hypertarget{sourcepointer-7}{}\fi
\noindent\textit{Source reference.} Original reference spaces-morphisms-lemma-composition-unramified: sources/stacks-algebraic-stacks-morphisms/pinned-original-source.tex, source lines 880--880. Navigation pointer only: the counted text refers to the Stacks item spaces-morphisms-lemma-composition-unramified, which belongs to another chapter of the project and is not part of this card. The pointer carries the target's own label as its number because this chapter's staged snapshot does not contain the other chapter's tag table; its content is not reproduced here.\par\medskip
\ifdefined\hypertarget\hypertarget{sourcepointer-8}{}\fi
\noindent\textit{Source reference.} Original reference stacks-properties-section-properties-morphisms: sources/stacks-algebraic-stacks-morphisms/pinned-original-source.tex, source lines 248--248. Navigation pointer only: the counted text refers to the Stacks item stacks-properties-section-properties-morphisms, which belongs to another chapter of the project and is not part of this card. The pointer carries the target's own label as its number because this chapter's staged snapshot does not contain the other chapter's tag table; its content is not reproduced here.\par\medskip
\begin{lemma}
\label{lemma-composition-separated}
All of the separation axioms listed in
Definition \ref{definition-separated}
are stable under composition of morphisms.
\end{lemma}

\begin{proof}
Let $f : \mathcal{X} \to \mathcal{Y}$ and
$g : \mathcal{Y} \to \mathcal{Z}$ be morphisms of algebraic stacks
to which the axiom in question applies.
The diagonal $\Delta_{\mathcal{X}/\mathcal{Z}}$ is the composition
$$
\mathcal{X} \longrightarrow
\mathcal{X} \times_\mathcal{Y} \mathcal{X} \longrightarrow
\mathcal{X} \times_\mathcal{Z} \mathcal{X}.
$$
Our separation axiom is defined by requiring the diagonal
to have some property $\mathcal{P}$. By
Lemma \ref{lemma-fibre-product-after-map}
above we see that the second arrow also has this property.
Hence the lemma follows since the composition of
morphisms which are representable by algebraic spaces with property
$\mathcal{P}$ also is a morphism with property $\mathcal{P}$, see
our general discussion in
Properties of Stacks,
Section \ref{stacks-properties-section-properties-morphisms}
and
Morphisms of Spaces, Lemmas
\ref{spaces-morphisms-lemma-composition-unramified},
\ref{spaces-morphisms-lemma-composition-quasi-finite},
\ref{spaces-morphisms-lemma-composition-proper},
\ref{spaces-morphisms-lemma-composition-quasi-compact}, and
\ref{spaces-morphisms-lemma-composition-separated}.
\end{proof}

\end{document}
